\section{中文输入与选择}
输入区域：原文。行内公式 $x^2+1$ 和项目宏 $\vect{x}$。

\begin{definition}{概率空间}{space}
样本空间 $\Omega$、事件与概率组成一个概率空间。
\end{definition}

\begin{theorem}{平方恒等式}{square}
\begin{equation}\label{eq:square}
(x+1)^2=x^2+2x+1.
\end{equation}
\end{theorem}
\begin{proof}
展开等式两侧即可得到结论。
\end{proof}

公式引用 \eqref{eq:square}，定理引用 \ref{thm:square}。

\includegraphics[width=.25\textwidth]{image.png}

\includegraphics[width=.25\textwidth]{sample.pdf}

\begin{tikzcd}
A \arrow[r,"f"] \arrow[d,"g"'] & B \arrow[d,"h"] \\
C \arrow[r,"k"'] & D
\end{tikzcd}

\begin{tabular}{cc}
中文 & 数值 \\
甲 & $1$ \\
乙 & $2$
\end{tabular}

\begin{align}
a &= b \\
\intertext{同一段文字}
c &= d
\end{align}

两个相同公式的编号应各自保留。

\begin{align}
a &= b \\
\intertext{同一段文字}
c &= d
\end{align}

拖选终点。测试到此结束。
